\documentclass[letterpaper,10.5pt]{article}

\usepackage{latexsym}
\usepackage[empty]{fullpage}
\usepackage{titlesec}
\usepackage{marvosym}
\usepackage[usenames,dvipsnames]{color}
\usepackage{verbatim}
\usepackage{enumitem}
\usepackage[hidelinks]{hyperref}
\usepackage{fancyhdr}
\usepackage[english]{babel}
\usepackage{tabularx}

%----------FONT OPTIONS----------
\usepackage[sfdefault]{FiraSans} 

\pagestyle{fancy}
\fancyhf{}
\fancyfoot{}
\renewcommand{\headrulewidth}{0pt}
\renewcommand{\footrulewidth}{0pt}

\addtolength{\oddsidemargin}{-0.5in}
\addtolength{\evensidemargin}{-0.5in}
\addtolength{\textwidth}{1in}
\addtolength{\topmargin}{-.5in}
\addtolength{\textheight}{1.0in}

\urlstyle{same}

\raggedbottom
\raggedright
\setlength{\tabcolsep}{0in}

\titleformat{\section}{
  \vspace{-4pt}\scshape\raggedright\large
}{}{0em}{}[\color{black}\titlerule \vspace{-5pt}]

\newcommand{\resumeItem}[1]{
  \item\small{
    {#1 \vspace{-2pt}}
  }
}

\newcommand{\resumeSubheading}[4]{
  \vspace{-2pt}\item
    \begin{tabular*}{0.97\textwidth}[t]{l@{\extracolsep{\fill}}r}
      \textbf{#1} & #2 \\
      \textit{\small#3} & \textit{\small #4} \\
    \end{tabular*}\vspace{-3pt}
}

\newcommand{\resumeProjectHeading}[2]{
    \item
    \begin{tabular*}{0.97\textwidth}{l@{\extracolsep{\fill}}r}
      \small#1 & #2 \\
    \end{tabular*}\vspace{-3pt}
}

\newcommand{\resumeSubItem}[1]{\resumeItem{#1}\vspace{-4pt}}

\renewcommand\labelitemii{$\vcenter{\hbox{\tiny$\bullet$}}$}

\newcommand{\resumeSubHeadingListStart}{\begin{itemize}[leftmargin=0.15in, label={}]}
\newcommand{\resumeSubHeadingListEnd}{\end{itemize}}
\newcommand{\resumeItemListStart}{\begin{itemize}}
\newcommand{\resumeItemListEnd}{\end{itemize}\vspace{-2pt}}

%-------------------------------------------

\begin{document}

%----------HEADING----------
\begin{center}
    {\Huge \scshape Ali Syed} \\ \vspace{2pt}
    \small College Station, TX $|$ 346-247-4146 $|$ \href{mailto:saabidi1@tamu.edu}{\underline{saabidi1@tamu.edu}} $|$
    \href{https://www.linkedin.com/in/ali-hamza-abidi-syed/}{\underline{linkedin.com/in/ali-syed}} $|$
    \href{https://scikitrom.github.io/}{\underline{Portfolio/Projects}}
\end{center}

%-----------SUMMARY-----------
\section*{Profile}
Full-stack ML engineer and applied mathematician who ships end-to-end---from data pipelines and distributed computing to production agentic AI services and customer-facing web applications. Proven ability to drive clarity and deliver outcomes in ambiguous, cross-functional environments.

%-----------EDUCATION-----------
\section{Education}
  \resumeSubHeadingListStart
    \resumeSubheading
      {Texas A\&M University}{College Station, TX}
      {Ph.D. in Applied Mathematics, GPA: 3.92}{Aug. 2022 -- Present}
      \resumeItemListStart
        \resumeItem{\textbf{Focus:} Deep learning, distributed computing, optimization, high-performance numerical methods.}
      \resumeItemListEnd

    \resumeSubheading
      {University of Houston}{Houston, TX}
      {B.Sc. in Mathematics, Minor in CS \& Data Science, GPA: 3.99}{Aug. 2019 -- May 2022}
  \resumeSubHeadingListEnd

%-----------EXPERIENCE-----------
\section{Work Experience}
  \resumeSubHeadingListStart

 \resumeSubheading
      {Texas A\&M Institute of Data Science (TAMIDS)}{College Station, TX}
      {ML Engineering \& Research Intern}{Jun. 2025 -- Aug. 2025}
      \resumeItemListStart
  \resumeItem{Designed and shipped \textbf{scikit-rom}, an open-source Python library for reduced-order modeling; drove requirements through weekly syncs with 4 domain-science research groups.}
  \resumeItem{Built a \textbf{PySpark} ETL pipeline to ingest, transform, and store large-scale FEM snapshot matrices in \textbf{Parquet/Delta Lake}, cutting data-prep time from hours to minutes.}
  \resumeItem{Deployed the full simulation stack as a \textbf{FastAPI} microservice with a \textbf{React} dashboard, enabling researchers to run ROM experiments from a browser.}
  \resumeItem{Containerized all services in \textbf{Docker} with \textbf{Jenkins} CI/CD; achieved \textbf{100--500x} speedup on linear solves and \textbf{$\sim$99\%} accuracy retention.}
\resumeItemListEnd

    \resumeSubheading
      {Switchless}{Remote / College Station, TX}
      {Machine Learning Engineer}{Jan. 2025 -- May 2025}
      \resumeItemListStart
  \resumeItem{Owned the full product lifecycle for an audio noise-detection system: ran customer discovery calls, defined milestones, and delivered iterative demos to converge on production requirements.}
  \resumeItem{Engineered a \textbf{PySpark + Delta Lake} pipeline to process \textbf{21 GB} of raw audio into versioned feature tables; automated feature extraction, validation, and \textbf{MLOps} lineage tracking.}
  \resumeItem{Built a production \textbf{Flask/REST API} with \textbf{Server-Sent Events (SSE)} streaming to push real-time inference results to a \textbf{React} monitoring dashboard used by the customer in the field.}
  \resumeItem{Deployed an Edge-ML CNN on Raspberry Pi via int8 quantization and pruning---\textbf{90\%} model-size reduction (45\,MB $\to$ 4.3\,MB), \textbf{90\%} accuracy retained.}
  \resumeItem{Reduced false positives by \textbf{35\%} through augmentation and targeted negative mining, directly resolving customer-reported field failures.}
\resumeItemListEnd
      
    \resumeSubheading
      {University of Houston}{Houston, TX}
      {Research Assistant -- Scientific Computing}{Jan. 2020 -- Aug. 2022}
      \resumeItemListStart
        \resumeItem{Built distributed numerical solvers in MATLAB and Julia for inverse acoustic problems on HPC clusters (\textbf{SLURM}, MPI); coordinated deliverables across physics and engineering teams.}
        \resumeItem{Developed chaos-analysis simulations (Lyapunov exponents, topological entropy) and presented reproducible results to faculty stakeholders, iterating on scope and methodology from feedback.}
      \resumeItemListEnd

  \resumeSubHeadingListEnd

%-----------PROJECTS-----------
\section{Technical Projects \& Research}
\resumeSubHeadingListStart

\resumeProjectHeading
{\textbf{Agentic Data Copilot} $|$ \emph{LangGraph, FastAPI, React, PySpark, Docker} $|$ \href{https://github.com/scikitrom/agentic-data-copilot}{\underline{GitHub}}}{Apr. 2025}
\resumeItemListStart
  \resumeItem{Built and deployed an \textbf{end-to-end agentic AI application}: users upload datasets, and a multi-agent pipeline autonomously plans, analyzes, generates code, executes it, and streams insights back in real-time.}
  \resumeItem{Architected a 5-node \textbf{LangGraph} agent graph (Planner $\to$ Analyst $\to$ Coder $\to$ Executor $\to$ Summarizer) with tool-use grounding---agents call pandas/plotly tools in a sandboxed runtime.}
  \resumeItem{Full-stack delivery: \textbf{FastAPI} backend with \textbf{RESTful} endpoints and \textbf{SSE} streaming, \textbf{React} frontend with drag-and-drop upload and live agent trace UI, \textbf{SQLite} session persistence.}
  \resumeItem{Automated data ingestion pipeline with profiling, feature extraction, and correlation analysis; packaged for one-command deployment via \textbf{Docker Compose}.}
\resumeItemListEnd

\resumeProjectHeading
{\textbf{DeepOSets -- In-Context Multi-Operator Learning} $|$ \emph{PyTorch, Distributed Training, W\&B}}{Jan. 2025 -- Present}
\resumeItemListStart
  \resumeItem{Formulated the in-context multi-operator learning problem and proved universality of a modified DeepOSets architecture.}
  \resumeItem{Scaled training across multi-GPU nodes via \textbf{PyTorch DDP} on TAMU HPC; managed experiment tracking and \textbf{MLOps} with \textbf{Weights \& Biases}.}
  \resumeItem{Achieved up to \textbf{84\%} MSE reduction across 12 PDE tasks, generalizing to unseen prompt sizes without fine-tuning.}
\resumeItemListEnd

\resumeProjectHeading
{\textbf{Group Equivariant Neural Networks} $|$ \emph{Python, PyTorch}}{Aug. 2024 -- Present}
\resumeItemListStart
  \resumeItem{Derived quantitative approximation rates for equivariant architectures (DeepSets, Transformers, frame-averaged models); published at \textit{arXiv}.}
  \resumeItem{Implemented frame-averaging layers in PyTorch, boosting MNIST detection from 25\% to \textbf{88\%} without changing the backbone.}
\resumeItemListEnd

\resumeSubHeadingListEnd

%-----------SKILLS-----------
\section{Technical Skills}
\begin{itemize}[leftmargin=0.15in, label={}]
\small{\item{
\textbf{Data Engineering:} Apache Spark, PySpark, Parquet, Delta Lake, ETL/ELT Pipelines, SQL, PostgreSQL \\
\textbf{Full-Stack:} FastAPI, Flask, RESTful APIs, Server-Sent Events (SSE), React, Docker Compose \\
\textbf{Agentic AI / ML:} LangChain, LangGraph, RAG, PyTorch, JAX, Scikit-learn, Transformers, CNNs, GNNs, Edge ML \\
\textbf{MLOps \& Infra:} Jenkins CI/CD, Docker, Git, Weights \& Biases, Linux, HPC (SLURM, MPI), PyTorch DDP \\
\textbf{Languages:} Python, C++, SQL, MATLAB, Julia, R
}}
\end{itemize}

%-----------PUBLICATIONS-----------
\section{Selected Publications \& Preprints}
  \resumeSubHeadingListStart
    \resumeItem{J. W. Siegel, S. Hordan, H. Lawrence, \textbf{A. Syed}, and N. Dym. ``Quantitative Approximation Rates for Group Equivariant Learning.'' \textit{arXiv:2602.20370}, 2026.}
    \resumeItem{S.-T. Chiu, A. Nambiar, \textbf{A. Syed}, J. W. Siegel, and U. Braga-Neto. ``In-Context Multi-Operator Learning with DeepOSets.'' \textit{arXiv:2512.16074}, 2025.}
  \resumeSubHeadingListEnd

\end{document}
